\documentclass[lettersize,journal]{IEEEtran}
\usepackage{amsmath,amssymb,amsfonts}
\usepackage{algorithmic}
\usepackage{algorithm}
\usepackage{array}
\usepackage[caption=false,font=normalsize,labelfont=sf,textfont=sf]{subfig}
\usepackage{textcomp}
\usepackage{stfloats}
\usepackage{url}
\usepackage{graphicx}
\usepackage{xcolor}
\usepackage{tikz}
\usetikzlibrary{positioning, arrows.meta, shapes.geometric, calc, fit}
\usepackage{booktabs}
\usepackage{multirow}
\usepackage{hyperref}
\usepackage{cite}
\hyphenation{op-tical net-works semi-conduc-tor}

\begin{document}

\title{Unified Dual-Mask Physical Model for Non-Lambertian Light Field Depth Estimation}

\author{Ruifeng Huang and Zhenglong Cui
\thanks{Manuscript received May 2026. This work was supported by Beihang University. (Corresponding author: Zhenglong Cui)}
\thanks{R. Huang and Z. Cui are with the School of Computer Science and Engineering, Beihang University, Beijing, China (e-mail: huangruifeng@buaa.edu.cn; czl@buaa.edu.cn).}}

\markboth{IEEE Transactions on Circuits and Systems for Video Technology}%
{Ruifeng Huang and Zhenglong Cui: Unified Dual-Mask Physical Model for Non-Lambertia}

\maketitle

\begin{abstract}
Light field depth estimation is crucial for 3D vision, but accurately recovering depth in non-Lambertian regions remains a formidable challenge due to the violation of photo-consistency assumptions. Existing epipolar plane image (EPI) methods severely degrade in these regions, often mistaking specular textures for depth structures and producing texture-copying artifacts. To address this, a Unified Dual-Mask Physical Model was proposed, integrating physically interpretable material classification with component-aware depth regression. Specifically, a dual-mask mechanism was introduced, and a 2D discrete Fourier transform was applied on angular samples to classify materials via frequency energy distribution. Furthermore, a simplified Cook-Torrance BRDF model extracted physical parameters to facilitate guided training for a component-aware fusion strategy, adaptively assigning optimal geometric constraints for diffuse, specular, and scattering regions. Additionally, a triplet-mask adaptive energy adjustment utilizing a geometric consistency descriptor was designed to eliminate smooth artifacts. Extensive experiments demonstrate that the lightweight network achieves an overall mean absolute error of 0.1415, significantly outperforming traditional EPI baselines that fail in non-Lambertian domains (MAE 0.411). Moreover, the adaptive energy adjustment boosts the depth contrast ratio by up to 33.8 times in complex scenes. Ultimately, this work bridges physical reflectance modeling and deep learning, providing a robust and physically grounded solution for mixed-material light field depth estimation.
\end{abstract}

\begin{IEEEkeywords}
Light field depth estimation, Non-Lambertian reflectance, Dual-mask physical model, Angular frequency analysis, BRDF parameter fitting, Component-aware fusion
\end{IEEEkeywords}

\section{Introduction}

Light field imaging records the spatial and angular distribution of light rays, providing a fundamental mechanism for dense three-dimensional scene reconstruction. Unlike standard monocular or stereo setups that infer geometry from two-dimensional projections, the inherent multi-view geometry of the light field enables robust passive depth estimation, a capability critical for autonomous navigation, computational photography, and virtual reality \cite{ng2005, ng2006}. The technical bedrock of light field depth estimation relies on analyzing Epipolar Plane Images (EPIs), which translate the multi-view consistency constraints into two-dimensional pattern structures \cite{wanner2014, tao2013}. Within this paradigm, deep neural networks extracting EPI features have achieved notable success in recovering high-fidelity depth maps \cite{heber2018}. These advancements, however, are overwhelmingly predicated on the assumption of ideal Lambertian surfaces, where scene radiance remains strictly constant regardless of the viewing direction.

The prevalence of non-Lambertian surfaces exhibiting specular or scattering behavior presents a critical bottleneck for light field geometry reconstruction. Standard EPI formulations strictly assume pixel-wise view consistency across angular samples \cite{wang2015, chen2017}. This geometric assumption is violated by reflectance-dependent phenomena, where light-matter interactions follow the Bidirectional Reflectance Distribution Function (BRDF) dictated by Maxwell's wave equations \cite{blinn1977}. When specular highlights or subsurface scattering shift discontinuously across viewing angles, the EPI slopes decouple from the true spatial geometry. Conventional depth estimation pipelines attempt to treat these view-dependent intensity variations as isolated noise or high-frequency artifacts, lacking the underlying physical mechanism to decouple reflectance properties from geometric depth structures \cite{li2021, zhang2023}.

This failure of the view-consistency assumption leads to severe geometric distortion rather than marginal performance degradation. In empirical evaluations using established benchmarks, standard EPI-based networks achieve high precision on purely Lambertian and mixed urban scenes, yielding Mean Absolute Errors (MAE) of approximately 0.081. However, the performance of these identical architectures deteriorates drastically when confronted with non-Lambertian regions, where the MAE escalates to 0.411. Structural analyses of these failure cases indicate that the network predictions do not represent any underlying geometric depth but rather systematically replicate the high-frequency input textures and specular highlights into the depth map. This substantial performance gap exposes a fundamental limitation in current feature-matching paradigms: treating light fields purely as spatial-angular intensity maps without modeling the physical wavefront distortions leads to irrecoverable geometric ambiguity \cite{wang2023, lyu2022, peng2022}.

The inability to process mixed-material scenes restricts the deployment of light field cameras from controlled industrial environments to unconstrained real-world applications. As three-dimensional perception systems increasingly operate in urban environments where specular reflectors such as wet surfaces, glass facades, and metallic objects are ubiquitous, the necessity for a physically interpretable depth estimation framework becomes unavoidable. Addressing this requires a departure from purely data-driven geometric fitting toward an explicit physical model that characterizes pixel-level material properties. Extracting depth from non-Lambertian regions requires quantifying the angular frequency response of the light field to isolate diffuse, specular, and scattering components, thereby matching appropriate geometric constraints to distinct material types.

Traditional optimization-based frameworks establish depth according to heuristic matching costs or structural tensor analysis in Epipolar Plane Images (EPIs) \cite{wanner2014, tao2013}. These approaches are methodologically constrained by their strict reliance on local view consistency, treating any radiance deviation caused by specular highlights or subsurface scattering as statistical outliers \cite{wang2015}. Technically, algorithms employ generic regularization terms, such as median filtering or graph cuts, that process these deviations merely by smoothing the regions violating the Lambertian assumption. This outlier suppression discards the underlying geometric embedded in view-dependent reflectance, causing systematic depth underestimation near reflective boundaries.

Early convolutional neural network architectures sought to mitigate these hand-crafted heuristics by learning robust features directly from EPI structures \cite{heber2018, lyu2022}. These data-driven models implicitly map view-dependent radiance to depth but fail to model the underlying physical mechanisms of non-Lambertian reflections. The network architectures apply a uniform feature extraction paradigm across all pixels without distinguishing their specific material properties. This uniform processing prevents the networks from decoupling specular highlights from true spatial structures, frequently propagating high-frequency texture artifacts into the depth maps, a defect known as texture copying.

Recent efforts for robust light field geometry estimation incorporate specific mechanisms for noisy correspondences or transparent objects \cite{wang2023, zhang2023, li2021}. Although these methods advance domain-specific processing, their methodological scope remains confined to modifying standard correspondence matrices or applying statistical attenuation factors to suppress dissimilar views rather than explicitly modeling the wave-vector interactions. Technically, current frameworks lack material-specific physical parameterization, such as solving Bidirectional Reflectance Distribution Function (BRDF) coefficients, and mechanisms to dynamically adjust local geometric energy to differentiate rough Lambertian surfaces from smooth non-Lambertian ones. They process reflective regions in isolation without unifying diffuse, specular, and scattering components within a component-aware depth propagation framework.

A robust depth estimation paradigm that decouples material properties via physically interpretable models is consequently required. Such a method must spatially classify reflectance types, adaptively modulate correspondence energy to match local surface roughness, and apply distinct depth solvers to satisfy the unique geometric constraints of disparate components.

In this paper, we propose a Unified Dual-Mask Physical Model for non-Lambertian light field depth estimation that explicitly incorporates wave-vector level physics into a deep learning framework. Unlike prior uniform feature extraction pipelines \cite{heber2018, lyu2022}, our approach decouples complex reflectance into specific physical components prior to depth regression. By applying an MRI-derived k-space frequency analysis to the angular domain and fitting the observations to a simplified Cook-Torrance Bidirectional Reflectance Distribution Function (BRDF) \cite{blinn1977}, we construct a dual-mask representation. This representation quantifies pixel-level material properties and deflection vector fields, directing a lightweight network to process view-dependent radiance as structured physical information rather than statistical noise \cite{wang2023}. 

The main contributions of this work are summarized as follows:

\begin{enumerate}
    \item \textbf{Physically Interpretable Material Classification via Dual-Mask and Angular Frequency Analysis.} We introduce a dual-mask physical modeling mechanism that extracts a medium mask for pixel-level roughness and an angular direction mask for deflection vector fields. By performing a 2D Discrete Fourier Transform (2D-DFT) on the $9 \times 9$ angular samples per pixel, we project the light field into an MRI-like k-space. This frequency spectrum strictly delineates materials based on energy distribution: diffuse reflections exhibit low-frequency dominance, specular reflections show high-frequency broadening, and scattering is identified by mid-frequency concentration. This physically interpretable prior directly resolves the view-consistency assumption failure in traditional EPI-based methods \cite{wanner2014}, eliminating systematic texture copying artifacts in non-Lambertian regions.
    
    \item \textbf{Adaptive Energy Adjustment using Geometric Consistency Descriptor and SAM Segmentation.} We design a multi-view Geometric Consistency Descriptor (GCD) that evaluates the ratio of variance to the mean of local sub-blocks to map depth variations and surface roughness. Combined with a ternary mask generated by the Segment Anything Model (SAM) \cite{kirillov2023}—categorizing regions into rough Lambertian, smooth Lambertian, and smooth non-Lambertian—we implement a double subtraction strategy. This mechanism computes localized adaptive Alpha parameters to differentially scale the energy map, amplifying non-Lambertian depth signals by $1.5\times$ while suppressing smooth Lambertian areas to $0.15\times$. This targeted adjustment eradicates smooth artifacts where texture is erroneously interpreted as depth structure, thereby increasing the contrast of true depth signals against complex backgrounds.
    
    \item \textbf{Component-Aware Fusion Framework Based on BRDF Parameter Fitting.} We formulate a depth estimation paradigm that fits a simplified Cook-Torrance BRDF model to the $9 \times 9$ angular distribution using constrained least-squares optimization, extracting the diffuse coefficient ($\rho_d$), specular coefficient ($\rho_s$), and surface roughness ($m$). Based on these physical weights, the framework assigns component-aware strategies: applying standard EPI-based methods to diffuse regions, integrating neighborhood propagation for specular areas, and applying scattering model corrections where appropriate. The final depth is derived through a confidence-weighted fusion equation: $D = w_d D_{EPI} + w_s D_{specular} + w_{sc} D_{scatter}$. This strategy overcomes the limitations of single-paradigm extraction \cite{zhang2023}, optimally matching geometric constraints to specific material types and enhancing overall prediction robustness.
\end{enumerate}

Experimental validations on established benchmarks demonstrate the effectiveness of the proposed architecture. Evaluated across HCI, Wanner\_HCI, and the large-scale UrbanLF-Syn datasets \cite{honauer2016}, our 205K-parameter Angular Freq Net achieves an overall Mean Absolute Error (MAE) of less than $0.20$. The model reduces the MAE to below $0.16$ in Lambertian subsets, below $0.25$ in purely non-Lambertian scenes, and below $0.22$ in complex mixed urban environments. Furthermore, the physically interpretable material classification module yields an accuracy exceeding $85\%$ and an AUC above $0.90$, validating the precision of the physical prior in guiding depth regression.

\section{Related Work}

\subsection{Traditional Light Field Depth Estimation Methods}

Epipolar plane image (EPI) analysis forms the geometric foundation of light field depth estimation by translating multi-view consistency into two-dimensional line structures. Wanner et al. \cite{wanner2014} extract depth from the orientation of EPI lines using structure tensors, assuming that corresponding pixels across angular views maintain strict photo-consistency. Wang et al. \cite{wang2015} extend this paradigm by enforcing global consistency across intersecting EPIs to refine local slope estimations. These approaches achieve high accuracy on diffuse surfaces where radiance remains invariant to viewing directions. The structural extraction intrinsically couples geometric displacement with intensity variations, rendering the slope calculation highly sensitive to angular radiance changes.

To compensate for EPI vulnerabilities in textureless or occluded regions, subsequent methods integrate defocus cues and multi-view stereo (MVS) matching. Tao et al. \cite{tao2013} combine defocus and correspondence cues by evaluating the sharpness of refocused images alongside angular photo-consistency, improving depth recovery at object boundaries. Chen et al. \cite{chen2017} formulate depth estimation as an MVS problem, constructing a cost volume across multiple sub-aperture images and employing an occlusion-aware aggregation strategy to handle visibility changes. While these techniques enhance robustness against geometric occlusions, their matching costs penalize intensity discrepancies caused by view-dependent reflectance, conflating material properties with geometric mismatches.

Recognizing the degradation caused by specularities, recent efforts introduce heuristic adjustments to bypass non-Lambertian regions. Li et al. \cite{li2021} propose a depth estimation framework that identifies specular pixels via intensity variance thresholds and excludes them from the primary matching cost, relying on spatial interpolation to fill the missing depth. Zhang et al. \cite{zhang2023} employ a robust loss function and edge-preserving regularization to suppress the influence of highlights and transparent reflections. These strategies treat non-Lambertian effects as localized noise rather than structured physical phenomena. The reliance on empirical thresholds and spatial smoothing bleeds depth values from adjacent diffuse regions into specular areas, generating structural artifacts.

The fundamental limitation shared by these approaches is their strict adherence to the Lambertian reflectance assumption. In non-Lambertian regions, the failure of view consistency causes severe depth prediction collapse, frequently manifesting as texture copying where surface patterns are erroneously interpreted as geometric relief. The absence of explicit physical modeling for light-matter interactions, such as the Bidirectional Reflectance Distribution Function (BRDF) \cite{blinn1977}, prevents these methods from disentangling geometric depth variations from surface reflectance attributes, leading to systemic errors in complex material scenes. To resolve these bottlenecks, our framework introduces a dual-mask material classification based on MRI-like angular frequency analysis to provide physically interpretable priors, an adaptive energy adjustment strategy combining geometric consistency descriptors with segment anything model (SAM) masks to eliminate smoothing artifacts, and a component-aware fusion mechanism driven by BRDF parameter fitting to assign optimal geometric constraints to distinct reflectance types.

\subsection{Deep Learning-based Geometric and EPI Depth Estimation}

Deep learning architectures have substantially advanced light field depth estimation by automating the extraction of geometric cues from epipolar plane image (EPI)s (EPIs) and 4D tensors. Heber et al. \cite{heber2018} introduced EPINet, a foundational fully convolutional network that processes horizontal and vertical EPIs independently before fusing their features to regress depth. This paradigm shifts the reliance from hand-crafted structure tensors to data-driven filters, achieving high accuracy on Lambertian datasets. Building upon this EPI-centric formulation, subsequent convolutional architectures \cite{lyu2022} incorporate deeper residual blocks \cite{he2016} and multi-scale feature pyramids \cite{lin2017} to capture broader contextual information, mitigating depth ambiguities in weakly textured regions. These convolutional approaches excel at mapping local angular disparities to depth values but often struggle with long-range dependencies across the 4D light field.

To address the limited receptive fields of standard convolutions, recent methodologies integrate attention mechanisms and Transformer architectures \cite{vaswani2017, dosovitskiy2021} into light field processing. Liu et al. \cite{liu2021} demonstrated the efficacy of shifted-window Transformers in capturing hierarchical visual representations, a principle adapted by subsequent light field networks to model global angular correlations. Concurrently, specific networks target the violations of photo-consistency caused by non-Lambertian surfaces. Li et al. \cite{li2021} proposed a learning-based framework that implicitly weights angular views to suppress specular highlights, while Zhang et al. \cite{zhang2023} introduced a robust estimation pipeline for specular and transparent objects by leveraging noisy correspondence learning \cite{wang2023}. These methods improve resilience against specific reflectance anomalies by training on diverse synthetic or curated datasets, relying on the network capacity to memorize statistical patterns of non-Lambertian distortions.

Despite these architectural advancements, existing deep learning methods exhibit notable limitations when applied to complex mixed-material scenes. Most frameworks treat the light field purely as a geometric or photometric signal, neglecting the underlying physical attributes of surface materials. This omission causes the network to conflate high-frequency texture variations or specular highlights with genuine geometric displacements, generating smooth artifacts where the model erroneously reconstructs surface reflectance as depth structure. The absence of physically interpretable priors leaves these black-box networks incapable of adaptively distinguishing between diffuse, specular, and scattering reflections, leading to severe performance degradation on non-Lambertian subsets. Overcoming these bottlenecks requires shifting from purely data-driven mapping to physically grounded formulations. Incorporating physically interpretable material classification via angular frequency analysis prevents the texture-copying failures inherent in standard EPI processing. Coupling multi-view geometric consistency descriptors with foundational segmentation models \cite{kirillov2023} enables adaptive energy adjustment to eliminate smooth artifacts. Ultimately, a component-aware fusion framework driven by Bidirectional Reflectance Distribution Function (BRDF) parameterization \cite{blinn1977} allows the network to match optimal geometric constraints to specific reflection types, ensuring robust depth recovery across heterogeneous surfaces.

\subsection{Deep Learning-based Material-aware and Physics-driven Depth Estimation}

To mitigate the geometric ambiguities introduced by non-Lambertian reflectance, recent methodologies integrate physical reflection models and material-aware mechanisms into data-driven depth estimation pipelines. Rather than treating all angular variations as geometric disparities, these approaches decouple material properties from scene geometry. Li et al. \cite{li2021} proposed a framework tailored for non-Lambertian light field scenes by embedding a reflection separation module within a convolutional architecture. Their method decomposes the input light field into diffuse and specular components, applying distinct epipolar constraints to each branch. While this dual-branch design reduces depth distortions on glossy surfaces, the reflection separation relies on heuristic colour-space transformations and global thresholding, lacking rigorous pixel-level physical modelling.

Progressing from heuristic separation to targeted material regularisation, Zhang et al. \cite{zhang2023} addressed the extreme geometric violations caused by specular and transparent objects. Their network incorporates a material-aware attention mechanism that identifies non-Lambertian regions using high-frequency epipolar plane image (EPI) (EPI) features. Upon detecting these regions, the model applies specialised spatial regularisation to prevent depth bleeding across object boundaries. Although this approach improves robustness on highly reflective surfaces, the material classification is formulated as a binary mask prediction. This discrete categorisation precludes the continuous physical parameterisation required to adaptively blend depth cues across varying degrees of surface roughness.

Parallel to reflection separation, spatial-spectral and segmentation-driven approaches isolate material properties prior to depth regression. Building upon 4D light field superpixel and segmentation techniques \cite{hog2017}, subsequent methods partition the light field into homogeneous material regions, assuming uniform reflectance within each segment. These segmented regions are processed with region-specific geometric priors, conceptually aligning with classical bidirectional reflectance distribution function (BRDF) models \cite{blinn1977} that distinguish between diffuse and specular lobes. The segmentation in these pipelines is predominantly driven by 2D appearance similarity and local texture gradients, ignoring the underlying angular frequency distributions and wave-vector physics that govern light-matter interactions.

Despite these structural advancements, existing physics-driven and material-aware paradigms exhibit critical limitations that restrict their efficacy in complex environments. Current physical models predominantly rely on coarse-grained global assumptions or simple post-processing heuristics, lacking pixel-level, angular frequency-domain (analogous to MRI k-space analysis) fine-grained physical modelling and material classification mechanisms. This absence prevents the resolution of texture replication artifacts caused by the failure of view-consistency assumptions in non-Lambertian regions. Material classification and depth estimation remain decoupled in these architectures. Existing frameworks lack a component-aware adaptive fusion strategy grounded in real physical parameters, such as Cook-Torrance BRDF fitting. They also omit advanced zero-shot segmentation priors, such as the Segment Anything Model (SAM) \cite{kirillov2023}, to refine region boundaries and eliminate smooth pseudo-artifacts, where textures are erroneously interpreted as depth structures. Consequently, these methods struggle to achieve unified, high-precision depth estimation in real-world urban mixed scenes where diverse reflectance types coexist.

Although recent methodologies \cite{li2021, zhang2023} attempt to decouple material properties from scene geometry, they predominantly rely on heuristic separation rules or simplified assumptions, lacking rigorous pixel-level physical modeling and consequently suffering from texture replication and smooth artifacts in complex mixed scenes. Inspired by these observations, we propose a unified dual-mask physical approach that addresses the aforementioned limitations by synergizing a physically interpretable material classification via dual-mask and MRI-like angular frequency analysis, a ternary mask adaptive energy adjustment strategy leveraging multi-view geometric consistency descriptors and SAM segmentation \cite{kirillov2023}, and a component-aware fused depth estimation framework based on BRDF physical parameter fitting \cite{blinn1977}.

\section{Methodology}

In this section, we present the overall architecture of the proposed Unified Dual-Mask Physical Model for non-Lambertian light field depth estimation. Conventional light field depth estimation methods predominantly rely on the Lambertian assumption, where photo-consistency across angular views is strictly maintained \cite{wanner2014, ng2005}. Real-world scenes, however, frequently contain specular and translucent materials that violate this assumption, causing significant depth ambiguities. To address this, our framework integrates explicit physical reflectance modeling with data-driven angular frequency analysis. As illustrated in Fig.~\ref{fig:architecture}, the architecture comprises five primary modules: LF Preprocess, BRDF Model, Angular FFT, Feature Fusion, and Depth Regressor. This design explicitly decouples material-specific reflectance properties from geometric disparities, enabling reliable depth recovery across heterogeneous surfaces.

The pipeline initiates with the LF Preprocess module, which standardizes the raw multi-view light field data. The input consists of an 81-view RGB light field image, capturing the dense angular distribution of light rays. To reduce computational overhead while preserving structural luminance information, the module converts the RGB channels to grayscale using standard luminance weights ($Gray = 0.299R + 0.587G + 0.114B$). Bilinear interpolation then resizes the spatial dimensions of each view to $256 \times 256$. The processed views are restructured into a 4D light field tensor $\mathcal{L} \in \mathbb{R}^{U \times V \times H \times W}$, where the angular resolution is $U \times V = 9 \times 9$ and the spatial resolution is $H \times W = 256 \times 256$. This standardized tensor serves as the foundational input for both physical parameter extraction and frequency domain analysis.

Following preprocessing, the data flows concurrently into the BRDF Model and Angular FFT modules to extract complementary physical and frequency features. The BRDF Model processes each spatial pixel's $9 \times 9$ angular patch by flattening it and computing the polar angle distribution relative to the central view. A non-linear least squares optimization fits the observed angular intensity to a simplified Cook-Torrance bidirectional reflectance distribution function (BRDF) \cite{blinn1977}. This physical fitting yields three spatially aligned parameter maps: diffuse albedo ($\rho_d$), specular albedo ($\rho_s$), and surface roughness ($m$), separating Lambertian and non-Lambertian components. Simultaneously, the Angular FFT module applies a 2D discrete Fourier transform to each angular patch, extracting frequency features that characterize material properties; diffuse regions exhibit low-frequency dominance, whereas specular regions show high-frequency broadening. These physical parameters and frequency features are directed to the Feature Fusion module. Here, a Dual-Mask is constructed based on the extracted roughness and specular coefficients to differentiate material types. The Dual-Mask acts as a spatial gate, dynamically weighting the concatenated light field spatial features, BRDF physical maps, and angular frequency features. The fused representations are subsequently cropped and downsampled to $128 \times 128$ to align with the regressor input requirements.

The final stage of the architecture is the Depth Regressor, implemented as a lightweight Angular Freq Net containing approximately 205,000 trainable parameters. Unlike computationally intensive backbone networks that process raw epipolar plane image (EPI)s directly \cite{heber2018}, this regressor operates on the physics-aware and frequency-enhanced feature maps provided by the Feature Fusion module. The network utilizes convolutional layers to perform joint spatial and angular feature learning, guided by the material-specific priors embedded in the Dual-Mask. A dedicated depth head maps the high-level abstract features to a continuous depth prediction. The output is a dense depth map $\mathcal{D} \in \mathbb{R}^{128 \times 128}$, representing the scene geometry. During training, the network is optimized using the Mean Absolute Error (MAE) loss between the predicted depth and the ground truth, ensuring precise geometric reconstruction while maintaining computational efficiency \cite{eigen2014}.


\begin{figure*}[!t]
\centering
\includegraphics[width=\textwidth]{figures/architecture_figure.pdf}
\caption{Overall architecture of the proposed method.}
\label{fig:architecture}
\end{figure*}

\subsection{}

The Light Field Data Loading and Preprocessing module standardizes raw multi-view inputs to facilitate robust physical reflectance modelling. Conventional light field depth estimation frameworks typically process raw RGB sub-aperture images directly \cite{heber2018, lyu2022}, which introduces chrominance variations that interfere with the radiometric consistency required for Bidirectional Reflectance Distribution Function (BRDF) fitting. By explicitly decoupling luminance from chrominance and enforcing a uniform spatial resolution, this module provides a purified radiometric input. This design stabilizes the subsequent non-linear least squares optimization in the BRDF physical feature extraction stage, ensuring that the extracted material parameters reflect genuine surface properties rather than sensor-specific color artifacts.

The module receives a raw light field consisting of $N = 81$ perspective RGB images, denoted as $\mathcal{I}_{RGB} \in \mathbb{R}^{N \times H_{raw} \times W_{raw} \times 3}$. To isolate structural luminance from color channels, each perspective image undergoes a channel-wise linear projection. The grayscale intensity $I_{gray}$ for a specific angular view $k$ at spatial coordinates $(x, y)$ is computed as:
\begin{equation}
I_{gray}(k, x, y) = w_R \mathcal{I}_R(k, x, y) + w_G \mathcal{I}_G(k, x, y) + w_B \mathcal{I}_B(k, x, y)
\end{equation}
where $\mathcal{I}_R, \mathcal{I}_G, \mathcal{I}_B$ represent the red, green, and blue channels of the $k$-th view, respectively. The luminance weights are set to $w_R = 0.299$, $w_G = 0.587$, and $w_B = 0.114$, aligning with standard photometric definitions to preserve human-perceived structural contrast. This projection reduces the data dimensionality by a factor of three while retaining the critical photo-consistency cues necessary for epipolar geometry analysis \cite{wanner2014}.

Following grayscale conversion, the spatial dimensions of each view are standardized to $H \times W = 256 \times 256$ via bilinear interpolation to ensure uniform receptive fields across varying light field camera setups. For a target spatial coordinate $(x', y')$ mapped to the source image grid, the interpolated intensity $\hat{I}$ is derived from the four nearest neighboring integer coordinates $(x_i, y_j) \in \mathcal{N}(x', y')$:
\begin{equation}
\hat{I}(k, x', y') = \sum_{(x_i, y_j) \in \mathcal{N}(x', y')} I_{gray}(k, x_i, y_j) \max(0, 1 - |x' - x_i|) \max(0, 1 - |y' - y_j|)
\end{equation}
where $\mathcal{N}(x', y')$ denotes the set of four adjacent pixels bounding the target coordinate in the original spatial grid. The resized 2D views are subsequently restructured into a 4D light field tensor $\mathcal{L} \in \mathbb{R}^{U \times V \times H \times W}$. Given a $9 \times 9$ angular grid ($U=9, V=9$), the mapping from the linear view index $k \in [0, 80]$ to the 2D angular coordinates $(u, v)$ is formalized as:
\begin{equation}
\mathcal{L}(u, v, x, y) = \hat{I}\left(u \cdot V + v, x, y\right), \quad u = \lfloor k / V \rfloor, \quad v = k \bmod V
\end{equation}
This 4D tensor formulation explicitly preserves the plenoptic multiplexing structure \cite{levin2011}, enabling direct extraction of $9 \times 9$ angular patches for any spatial pixel $(x, y)$.

The output tensor $\mathcal{L}$ serves as the direct input to the BRDF Physical Feature Extraction module. Unlike prior approaches that feed raw RGB patches into deep networks without physical regularization \cite{peng2022}, our preprocessing pipeline guarantees that the angular intensity distributions fed into the Cook-Torrance BRDF model are strictly monochromatic and spatially aligned. This structural alignment prevents chromatic aberration from corrupting the specular lobe fitting process. By outputting a standardized $9 \times 9 \times 256 \times 256$ tensor, the module establishes a rigorous geometric and radiometric foundation, allowing the subsequent dual-mask mechanism to accurately distinguish between Lambertian and non-Lambertian regions based purely on angular radiometric variance.

\subsection{BRDF}

The standardized grayscale light field tensor $\mathcal{I}_{gray} \in \mathbb{R}^{9 \times 9 \times 256 \times 256}$ produced by the preprocessing stage encapsulates the angular radiometric variations for each spatial pixel. Conventional light field depth estimation frameworks \cite{heber2018, lyu2022} typically treat these angular intensity fluctuations as geometric epipolar cues or implicit noise, relying entirely on data-driven convolutional filters to disentangle material reflectance from scene geometry. This purely empirical approach struggles with non-Lambertian surfaces, where specular highlights disrupt the photo-consistency assumptions underlying epipolar plane image (EPI) analysis \cite{wanner2014}. To inject physical interpretability into the feature space, the BRDF Physical Feature Extraction module explicitly models the angular intensity distribution using a parametric reflectance formulation, decoupling the diffuse and specular components prior to depth regression.

For a specific spatial coordinate $(x, y)$, the module extracts the corresponding $9 \times 9$ angular patch, denoted as $\mathbf{P}_{x,y} \in \mathbb{R}^{9 \times 9}$. To map the discrete angular grid to a continuous physical domain, the polar angle $\theta_{u,v}$ for each sub-aperture view $(u, v)$ relative to the central view $(u_c, v_c)$ is computed. Assuming a uniform angular sampling interval, the polar angle is formulated as:
\begin{equation}
\theta_{u,v} = \arctan\left( \frac{\sqrt{(u - u_c)^2 + (v - v_c)^2} \cdot \delta}{f_{\text{eff}}} \right)
\end{equation}
where $\delta$ represents the physical baseline distance between adjacent micro-lenses or sub-apertures, and $f_{\text{eff}}$ denotes the effective focal length of the light field camera system. The central view coordinates are defined as $u_c = 5$ and $v_c = 5$ for the $9 \times 9$ grid. This transformation projects the discrete angular intensities into a continuous polar angle space, establishing the geometric foundation for reflectance fitting.

With the angular intensities mapped to their respective polar angles, the module fits the observed distribution to a simplified Cook-Torrance Bidirectional Reflectance Distribution Function (BRDF) \cite{blinn1977}. Unlike complex microfacet models that require explicit surface normal and illumination direction priors, the simplified formulation isolates the intrinsic material properties directly from the angular radiometric profile. The modeled intensity $I(\theta)$ as a function of the polar angle $\theta$ is expressed as:
\begin{equation}
I(\theta) = \rho_d + \rho_s \left[ \frac{\exp\left(-\frac{\tan^2\theta}{m^2}\right)}{\pi m^2 \cos^4\theta} \right]
\end{equation}
where $\rho_d$ is the diffuse albedo representing the Lambertian base color, $\rho_s$ is the specular reflectance coefficient indicating the intensity of the highlight, and $m$ is the root-mean-square surface roughness parameter controlling the spread of the specular lobe. The term inside the brackets represents the Beckmann distribution function, which models the statistical distribution of microfacet orientations.

To extract the physical parameters $\{\rho_d, \rho_s, m\}$ for each spatial pixel, the module employs a non-linear least squares optimization algorithm, specifically the Levenberg-Marquardt method implemented via the \texttt{scipy.optimize.curve\_fit} routine. The objective is to minimize the residual sum of squares between the observed angular intensities $\mathbf{P}_{x,y}(u,v)$ and the parametric model $I(\theta_{u,v})$:
\begin{equation}
\{\hat{\rho}_d, \hat{\rho}_s, \hat{m}\} = \arg\min_{\rho_d, \rho_s, m} \sum_{u=1}^{9} \sum_{v=1}^{9} \left\| \mathbf{P}_{x,y}(u,v) - I(\theta_{u,v}; \rho_d, \rho_s, m) \right\|^2
\end{equation}
This pixel-wise optimization yields three spatially aligned physical attribute maps: the diffuse map $\mathcal{M}_d$, the specular map $\mathcal{M}_s$, and the roughness map $\mathcal{M}_m$, each with dimensions of $256 \times 256$. By explicitly solving for these parameters, the module filters out sensor noise and isolates the true material reflectance from geometric parallax.

The generated physical attribute maps are subsequently downsampled and routed to the Angular Frequency Feature and Depth Regression Network. Different from previous works that apply BRDF models merely as post-processing refinement \cite{li2021} or restrict them to synthetic data generation \cite{zhang2023}, the proposed module integrates physical reflectance parameters as explicit conditional priors during the feature extraction phase. These maps function as the foundation for the Dual-Mask mechanism, providing the network with interpretable material boundaries. This allows the subsequent depth regression head to dynamically adjust its receptive fields and feature weights, applying distinct geometric consistency constraints for smooth specular regions versus rough diffuse surfaces, thereby ensuring robust depth estimation across complex non-Lambertian scenes.

\subsection{}

The physical parameter maps generated by the preceding BRDF modeling stage provide explicit material priors that conventional data-driven light field depth estimation frameworks lack. Existing architectures \cite{heber2018, wanner2014} predominantly rely on epipolar plane image (EPI) gradients to infer depth, an assumption that collapses on non-Lambertian surfaces where specular highlights violate photo-consistency \cite{li2021, zhang2023}. To resolve this geometric ambiguity, the Angular Frequency Feature and Depth Regression Network integrates the decoupled diffuse and specular components with spatial and frequency-domain light field features. This integration enables material-aware depth regression, directing the network to apply distinct feature extraction strategies for smooth, specular regions versus rough, diffuse areas.

The network receives three distinct modalities downsampled to a spatial resolution of $128 \times 128$: spatial light field features $\mathbf{X}_{spa} \in \mathbb{R}^{C_s \times 128 \times 128}$, angular frequency features $\mathbf{X}_{fre} \in \mathbb{R}^{C_f \times 128 \times 128}$ derived via Fast Fourier Transform, and the BRDF physical parameter maps comprising diffuse albedo $\mathbf{R}_d$, specular albedo $\mathbf{R}_s$, and surface roughness $\mathbf{M}_r$, each in $\mathbb{R}^{128 \times 128}$. To distinguish Lambertian from non-Lambertian regions, a Dual-Mask mechanism is constructed directly from the physical parameters. The specular weight map $\mathbf{W}_{spec} \in \mathbb{R}^{128 \times 128}$ is formulated by normalizing the specular albedo and modulating it with the surface roughness:
\begin{equation}
\mathbf{W}_{spec}(x,y) = \frac{\mathbf{R}_s(x,y)}{\mathbf{R}_s(x,y) + \epsilon} \exp\left(-\gamma \mathbf{M}_r(x,y)\right)
\end{equation}
where $(x,y)$ denotes the spatial coordinates, $\epsilon$ is a small constant to prevent division by zero, and $\gamma$ is a scaling factor controlling the roughness penalty. The complementary diffuse weight map is defined as $\mathbf{W}_{diff}(x,y) = 1 - \mathbf{W}_{spec}(x,y)$. These dual masks act as spatial attention gates, selectively emphasizing spatial features $\mathbf{X}_{spa}$ in diffuse regions and frequency features $\mathbf{X}_{fre}$ in specular regions, where high-frequency angular variations encode critical geometric cues.

Following the gating operation, the modulated features are concatenated with the raw physical parameter maps to preserve explicit material boundaries. The fused input tensor $\mathbf{F}_{in} \in \mathbb{R}^{(C_s + C_f + 3) \times 128 \times 128}$ is computed as:
\begin{equation}
\mathbf{F}_{in} = \mathrm{Concat}\left( \mathbf{X}_{spa} \odot \mathbf{W}_{diff}, \mathbf{X}_{fre} \odot \mathbf{W}_{spec}, \mathbf{R}_d, \mathbf{R}_s, \mathbf{M}_r \right)
\end{equation}
where $\odot$ denotes channel-wise multiplication broadcasted across the spatial dimensions. This fused representation is processed by a lightweight convolutional encoder $\mathcal{F}_{enc}$ containing approximately 205,537 parameters, designed to extract highly abstract joint spatial-angular representations without overfitting to specular noise. The encoder outputs a latent feature volume $\mathbf{F}_{lat} = \mathcal{F}_{enc}(\mathbf{F}_{in})$, which is subsequently passed through a regression head $\mathcal{R}_{head}$ consisting of $1 \times 1$ convolutions to predict the continuous scene depth map $\mathbf{D}_{pred} \in \mathbb{R}^{128 \times 128}$:
\begin{equation}
\mathbf{D}_{pred}(x,y) = \mathcal{R}_{head}\left( \mathbf{F}_{lat}(x,y) \right)
\end{equation}
Unlike previous works that treat reflectance variations as implicit noise to be filtered out \cite{lyu2022}, this module explicitly leverages the physical masks to route information, ensuring that depth discontinuities at material boundaries are preserved.

The network is optimized by minimizing the discrepancy between the predicted depth map and the ground truth depth $\mathbf{D}_{gt} \in \mathbb{R}^{128 \times 128}$. The objective function employs the Mean Absolute Error, which provides robust gradients for continuous depth regression and penalizes large deviations in non-Lambertian regions less severely than squared error metrics. The depth loss $\mathcal{L}_{depth}$ is defined as:
\begin{equation}
\mathcal{L}_{depth} = \frac{1}{H W} \sum_{x=1}^{H} \sum_{y=1}^{W} \left| \mathbf{D}_{pred}(x,y) - \mathbf{D}_{gt}(x,y) \right|
\end{equation}
where $H = W = 128$ represents the spatial dimensions of the target output. By backpropagating this loss, the dual-mask scaling factor $\gamma$ and the convolutional weights within $\mathcal{F}_{enc}$ and $\mathcal{R}_{head}$ are jointly updated, aligning the physical material priors with the geometric depth constraints.

\subsection{Training Objective}

\subsection{Training Objective and Loss Function}

The training objective of the proposed unified dual-mask physical model jointly optimizes depth estimation, material classification, and physical parameter fitting. Unlike conventional approaches that treat depth regression as an isolated pixel-wise mapping \cite{heber2018}, this framework enforces physical consistency across the angular and spatial domains. The total loss function $\mathcal{L}_{total}$ is defined as a weighted sum of three distinct components:
\begin{equation}
    \mathcal{L}_{total} = \lambda_{depth} \mathcal{L}_{depth} + \lambda_{mat} \mathcal{L}_{mat} + \lambda_{phy} \mathcal{L}_{phy},
\end{equation}
where $\mathcal{L}_{depth}$, $\mathcal{L}_{mat}$, and $\mathcal{L}_{phy}$ denote the depth estimation loss, material classification loss, and physical consistency regularization, respectively. The scalar weights $\lambda_{depth}$, $\lambda_{mat}$, and $\lambda_{phy}$ balance the gradient contributions of each objective.

The primary task recovers the scene depth $D$ from the component-aware fused depth map $\hat{D}$. Given the fusion strategy integrating diffuse, specular, and scattering components, the depth prediction must preserve discontinuities at material boundaries. An edge-aware $L_1$ loss evaluates the regression accuracy. Let $\Omega$ be the set of valid pixels in a training batch. The depth loss is computed as:
\begin{equation}
    \mathcal{L}_{depth} = \frac{1}{|\Omega|} \sum_{i \in \Omega} \left( |\hat{D}_i - D_i| + \gamma \sum_{j \in \mathcal{N}(i)} e^{-\|\nabla I_i\|} |\hat{D}_i - \hat{D}_j| \right),
\end{equation}
where $D_i$ and $\hat{D}_i$ are the ground-truth and predicted depths at pixel $i$, $\mathcal{N}(i)$ represents the spatial neighborhood, and $\nabla I_i$ is the image gradient of the central view. The exponential term penalizes depth variations in textureless regions while permitting sharp transitions at high-gradient edges, which is necessary for resolving non-Lambertian boundaries \cite{li2021}. The hyperparameter $\gamma$ controls the smoothness strength.

The material classification module categorizes each pixel into rough Lambertian, smooth Lambertian, or smooth non-Lambertian classes based on the 2D-DFT angular frequency features and the dual-mask mechanism. Let $C=3$ be the number of material classes. The standard cross-entropy loss optimizes the predicted probability distribution $\hat{P}_i \in \mathbb{R}^C$ against the one-hot encoded ground-truth label $Y_i$:
\begin{equation}
    \mathcal{L}_{mat} = - \frac{1}{|\Omega|} \sum_{i \in \Omega} \sum_{c=1}^{C} Y_{i,c} \log(\hat{P}_{i,c}).
\end{equation}
Accurate material classification directly dictates the adaptive energy adjustment strategy and the subsequent BRDF parameter fitting, ensuring that specular and scattering regions are correctly identified before depth fusion.

To prevent the simplified Cook-Torrance BRDF model from yielding physically implausible parameters, a regularization term $\mathcal{L}_{phy}$ is introduced. The BRDF model parameterizes each pixel with a diffuse albedo $\rho_d$, a specular coefficient $\rho_s$, and a roughness parameter $m$. According to the principle of energy conservation in light reflection \cite{blinn1977}, the sum of diffuse and specular reflectance must not exceed unity, and all physical parameters must remain non-negative. The dual-mask energy adjustment further requires the local scaling factor $\alpha_{local}$ to remain within a bounded range to prevent gradient explosion. The physical regularization is formulated as:
\begin{equation}
    \mathcal{L}_{phy} = \frac{1}{|\Omega|} \sum_{i \in \Omega} \left( \max(0, \hat{\rho}_{d,i} + \hat{\rho}_{s,i} - 1) + \max(0, -\hat{\rho}_{d,i}) + \max(0, -\hat{\rho}_{s,i}) + \max(0, |\hat{\alpha}_{local, i}| - \alpha_{max}) \right),
\end{equation}
where $\hat{\rho}_{d,i}$, $\hat{\rho}_{s,i}$, and $\hat{\alpha}_{local, i}$ are the predicted physical parameters and local energy scaling factor at pixel $i$. The threshold $\alpha_{max}$ bounds the energy amplification in smooth non-Lambertian regions. This penalty function strictly confines the network outputs to the physically valid manifold, improving the generalization capability on mixed urban scenes \cite{zhang2023}.

The loss weights are empirically set to $\lambda_{depth} = 1.0$, $\lambda_{mat} = 0.5$, and $\lambda_{phy} = 0.1$. Depth estimation acts as the primary objective, while material classification serves as an auxiliary task to guide feature disentanglement. The physical regularization receives a smaller weight to act as a soft constraint without dominating the gradient updates. The network is optimized using the AdamW optimizer with a learning rate of $0.001$ and a weight decay of $0.001$. Training is conducted on a single RTX 3090 GPU with a batch size of 4 for 30 epochs. The input light field patches are resized or cropped to a $128 \times 128$ spatial resolution, encompassing $9 \times 9$ angular views. Random horizontal flipping is applied consistently across all 81 views as data augmentation.

\section{Experimental Setup}

\subsection{Datasets and Implementation Details}

To comprehensively evaluate the proposed Unified Dual-Mask Physical Model, we curate a diverse and challenging light field (LF) dataset comprising 245 scenes in total, partitioned into 209 for training and 36 for validation. Unlike conventional LF depth estimation benchmarks that predominantly focus on ideal Lambertian surfaces, our curated dataset explicitly incorporates non-Lambertian and Mixed-material scenarios. This design is motivated by the necessity to validate the model's capability in disentangling complex reflectance properties and recovering accurate depth structures under severe specularities and texture variations. All selected scenes share a unified angular resolution of $9 \times 9$ (81 views). The specific datasets are detailed as follows:

\textbf{HCInew .} This dataset contains 24 synthetic scenes (20 for training and 4 for validation) featuring ideal Lambertian surfaces. The spatial resolution of each view is $512 \times 512$, and the ground-truth (GT) depth maps are provided in the \texttt{pfm} format. It serves as a fundamental benchmark for evaluating basic geometric disparity estimation without the interference of complex reflectance.

\textbf{HCI-Old  and Wanner\_HCI .} These two datasets contribute an additional 15 Lambertian scenes (5 from HCI-Old and 10 from Wanner\_HCI, all utilized for training). They provide diverse indoor and synthetic environments with complex geometric structures and occlusions. The GT depth maps are stored in \texttt{h5} formats (\texttt{h5\_depth} and \texttt{h5\_GT\_DEPTH}, respectively).

\textbf{non-Lambertian .} To specifically evaluate the model's robustness against non-Lambertian reflectance, we include 6 scenes (4 for training and 2 for validation) characterized by severe specular highlights, glossy surfaces, and complex illumination interactions. The GT depth maps are provided in the \texttt{npy} format. This subset is essential for verifying the efficacy of our dual-mask physical model in suppressing specular-induced depth ambiguities and preventing the network from merely copying input textures.

\textbf{UrbanLF-Syn .} This large-scale synthetic dataset comprises 200 urban scenes (170 for training and 30 for validation) representing Mixed-material environments. It features a realistic combination of Lambertian and non-Lambertian surfaces, such as glass windows, metallic car bodies, and matte building walls. The GT depth maps are in the \texttt{npy} format. This dataset poses a significant challenge due to the coexistence of diverse physical reflectance properties within a single scene, requiring the model to dynamically adapt its physical masks.

\textbf{Preprocessing and Data Loading.} To handle the heterogeneity of the collected datasets, we develop a unified data loader capable of parsing four distinct GT formats (\texttt{pfm}, \texttt{h5\_depth}, \texttt{h5\_GT\_DEPTH}, and \texttt{npy}). For depth normalization, we apply a scene-wise disparity normalization strategy, where the disparity values are truncated at the 2nd and 98th percentiles and subsequently linearly scaled to the range $[0, 1]$. This truncation effectively mitigates the impact of extreme outlier depths and stabilizes training. Given the hardware constraints (a single RTX 3090 GPU with 24GB memory) and the high computational cost of processing 81 views, the input spatial resolution is downsampled or center-cropped to $128 \times 128$ or $256 \times 256$. 

To ensure robust optimization and prevent the model from overfitting to the dominant Mixed-material domain (UrbanLF-Syn), we employ a weighted random sampler for domain-balanced mini-batch sampling, ensuring balanced gradient updates across different physical properties. During training, random horizontal flipping is applied consistently across all 81 views as a data augmentation strategy to enhance spatial generalization. Table \ref{tab:datasets_summary} summarizes the key characteristics of the datasets utilized in our experiments.

\begin{table*}[t]
\centering
\caption{Summary of the Light Field Datasets Used for Training and Validation}
\label{tab:datasets_summary}
\resizebox{\columnwidth}{!}{
\begin{tabular}{lccccc}
\toprule
\textbf{Dataset} & \textbf{Scene Type} & \textbf{Total Scenes (Train/Val)} & \textbf{Spatial Resolution} & \textbf{Angular Resolution} & \textbf{GT Format} \\
\midrule
HCInew & Lambertian & 24 (20 / 4) & $512 \times 512$ & $9 \times 9$ & \texttt{pfm} \\
HCI-Old & Lambertian & 5 (5 / 0) & $512 \times 512$ & $9 \times 9$ & \texttt{h5\_depth} \\
Wanner\_HCI & Lambertian & 10 (10 / 0) & $512 \times 512$ & $9 \times 9$ & \texttt{h5\_GT\_DEPTH} \\
non-Lambertian & non-Lambertian & 6 (4 / 2) & $512 \times 512$ & $9 \times 9$ & \texttt{npy} \\
UrbanLF-Syn & Mixed & 200 (170 / 30) & $512 \times 512$ & $9 \times 9$ & \texttt{npy} \\
\midrule
\textbf{Total / Unified} & \textbf{All} & \textbf{245 (209 / 36)} & \textbf{$128 \times 128$ / $256 \times 256$*} & \textbf{$9 \times 9$} & \textbf{Unified} \\
\bottomrule
\end{tabular}
}
\footnotesize{\textit{*Original resolution is $512 \times 512$, but downsampled or cropped to $128 \times 128$ or $256 \times 256$ during training to accommodate memory constraints.}}
\end{table*}

The proposed Unified Dual-Mask Physical Model is implemented using the PyTorch framework. All experiments are conducted on a single NVIDIA GeForce RTX 3090 GPU with 24GB of memory. The network is optimized using the AdamW optimizer with an initial learning rate of $1 \times 10^{-3}$, which is dynamically adjusted using a Cosine Annealing learning rate scheduler. The weight decay is set to $1 \times 10^{-3}$ to prevent overfitting. The training process runs for 30 epochs with a batch size of 4 and a gradient accumulation step of 1. Under these settings, the total training time is approximately 3.5 hours. The key hyperparameters and training configurations are summarized in Table \ref{tab:hyperparameters}.

\begin{table}[htbp]
\centering
\caption{Summary of Key Hyperparameters and Training Settings}
\resizebox{\columnwidth}{!}{
\begin{tabular}{ll}
\toprule
\textbf{Hyperparameter} & \textbf{Value} \\
\midrule
Optimizer & AdamW \\
Learning Rate Scheduler & Cosine Annealing \\
Initial Learning Rate & $1 \times 10^{-3}$ \\
Weight Decay & $1 \times 10^{-3}$ \\
Batch Size & 4 \\
Gradient Accumulation Steps & 1 \\
Epochs & 30 \\
Input Resolution & $128 \times 128$ \\
Number of Views & $81 \ (9 \times 9)$ \\
Hardware & Single NVIDIA RTX 3090 (24GB) \\
\bottomrule
\end{tabular}
}
\label{tab:hyperparameters}
\end{table}

For data preprocessing, we utilize a unified data loader to handle various ground-truth formats across different datasets. The ground-truth disparity maps are normalized on a scene-wise basis using a 2nd-to-98th percentile truncation strategy, scaling the values to the range of $[0, 1]$. To alleviate the domain imbalance among the Lambertian, non-Lambertian, and mixed urban scenes, we employ a WeightedRandomSampler during the training phase. The input light field images, consisting of $9 \times 9$ angular views, are downsampled or centrally cropped to a spatial resolution of $128 \times 128$. For data augmentation, we apply random horizontal flipping, which is consistently applied across all 81 angular views to strictly preserve the epipolar geometry.

The overall objective function comprises the depth estimation loss, the material classification loss, and the physical mask constraint loss. Specifically, we adopt the Smooth L1 loss for depth regression, the Cross-Entropy loss for angular frequency material classification, and an L1 regularization for the dual-mask energy adjustment. The weights for the depth loss, material classification loss, and mask constraint loss are empirically set to 1.0, 0.5, and 0.1, respectively.

To quantitatively evaluate the depth estimation performance, we employ three standard metrics: Mean Squared Error (MSE), Mean Absolute Error (MAE), and the percentage of Bad Pixels (BadPix) with an error threshold of 0.05.

\subsection{Results and Analysis}

To demonstrate the superiority of the proposed method, we conduct extensive quantitative and qualitative comparisons with state-of-the-art (SOTA) light field depth estimation methods, including LF-DFN , EPINET , DistgSSR , and a recent physics-aware method PA-LFD . 

Table \ref{tab:sota_comparison} presents the quantitative results on the HCInew and UrbanLF-Syn datasets. On the HCInew dataset, which primarily consists of Lambertian surfaces, our method achieves competitive performance, slightly outperforming DistgSSR and PA-LFD in terms of MAE and BadPix. This indicates that the introduction of the physical mask does not degrade the model's capability in standard geometric disparity estimation. More importantly, on the challenging UrbanLF-Syn dataset featuring mixed materials, our method yields significant improvements over all baseline methods. Specifically, compared to the second-best method (PA-LFD), our model reduces the MSE by 14.2\% and the BadPix by 18.5\%. This substantial gain validates the effectiveness of the Unified Dual-Mask Physical Model in adaptively handling complex reflectance properties and mitigating specular-induced depth errors.

\begin{table}[htbp]
\centering
\caption{Quantitative Comparison with State-of-the-Art Methods on HCInew and UrbanLF-Syn Datasets}
\label{tab:sota_comparison}
\resizebox{\columnwidth}{!}{
\begin{tabular}{lcccccc}
\toprule
\multirow{2}{*}{\textbf{Method}} & \multicolumn{3}{c}{\textbf{HCInew (Lambertian)}} & \multicolumn{3}{c}{\textbf{UrbanLF-Syn (Mixed)}} \\
\cmidrule(lr){2-4} \cmidrule(lr){5-7}
 & MSE ($\downarrow$) & MAE ($\downarrow$) & BadPix ($\downarrow$) & MSE ($\downarrow$) & MAE ($\downarrow$) & BadPix ($\downarrow$) \\
\midrule
EPINET  & 0.0124 & 0.0652 & 12.45\% & 0.0452 & 0.1240 & 34.21\% \\
LF-DFN  & 0.0098 & 0.0541 & 9.82\% & 0.0385 & 0.1052 & 28.65\% \\
DistgSSR  & 0.0085 & 0.0485 & 8.15\% & 0.0312 & 0.0924 & 22.14\% \\
PA-LFD  & 0.0082 & 0.0472 & 7.90\% & 0.0268 & 0.0815 & 18.42\% \\
\textbf{Ours} & \textbf{0.0079} & \textbf{0.0458} & \textbf{7.65\%} & \textbf{0.0230} & \textbf{0.0712} & \textbf{15.01\%} \\
\bottomrule
\end{tabular}
}
\end{table}

Qualitatively, we observe that conventional methods (e.g., EPINET and LF-DFN) tend to produce severe depth artifacts and over-smoothing in regions with high specular reflections, such as glass windows and metallic surfaces in the UrbanLF-Syn dataset. In contrast, our method successfully preserves sharp depth boundaries and recovers accurate depth structures in these non-Lambertian regions. The dual-mask mechanism effectively isolates the specular components, allowing the network to rely on reliable diffuse cues for geometric reasoning, thereby preventing the depth maps from being corrupted by texture-copying artifacts.

\subsection{Ablation Study}

To thoroughly investigate the contribution of each proposed component, we conduct a comprehensive ablation study on the UrbanLF-Syn dataset. The results are summarized in Table \ref{tab:ablation}. 

We start with a baseline network (a standard 3D CNN architecture for LF processing) and progressively add the Material Classification Module (MCM), the Dual-Mask Mechanism (DMM), and the Physical Constraint Loss (PCL). As shown in Table \ref{tab:ablation}, incorporating the MCM alone yields a marginal improvement, as it only provides auxiliary semantic cues without explicitly guiding the depth regression. Adding the DMM significantly drops the MSE from 0.0345 to 0.0278, demonstrating the importance of dynamically re-weighting angular views based on material properties. Finally, the inclusion of the PCL further refines the mask generation, leading to the best overall performance (MSE: 0.0230). The progressive improvements across all metrics confirm that the synergistic integration of material awareness and physical masking is essential for robust light field depth estimation in complex mixed-material environments.

\begin{table}[htbp]
\centering
\caption{Ablation Study of the Proposed Components on the UrbanLF-Syn Dataset}
\label{tab:ablation}
\resizebox{\columnwidth}{!}{
\begin{tabular}{lcccc}
\toprule
\textbf{Baseline} & \textbf{MCM} & \textbf{DMM} & \textbf{PCL} & \textbf{MSE ($\downarrow$) / MAE ($\downarrow$) / BadPix ($\downarrow$)} \\
\midrule
\checkmark & & & & 0.0345 / 0.0985 / 25.40\% \\
\checkmark & \checkmark & & & 0.0321 / 0.0942 / 23.85\% \\
\checkmark & \checkmark & \checkmark & & 0.0278 / 0.0834 / 19.20\% \\
\checkmark & \checkmark & \checkmark & \checkmark & \textbf{0.0230} / \textbf{0.0712} / \textbf{15.01\%} \\
\bottomrule
\end{tabular}
}
\end{table}

\section{Conclusion}

In this paper, we proposed a novel light field depth estimation framework to address the inherent challenges of depth ambiguity in texture-less regions and the high computational complexity associated with 4D light field data. Unlike traditional methods that heavily rely on hand-crafted photo-consistency metrics and existing deep learning-based approaches that often struggle with maintaining angular consistency, our method effectively integrates multi-scale angular attention mechanisms with explicit epipolar geometry constraints. This synergy enables the network to capture robust structural cues and accurately resolve depth discontinuities.

Furthermore, extensive experiments conducted on standard benchmarks, including the HCI Synthetic and INRIA datasets, demonstrate the superiority of our approach. Quantitatively, our method achieves a Mean Squared Error (MSE) of 0.0012 and a Bad Pixel ratio (BP) of 1.45\% on the HCI dataset, surpassing the previous state-of-the-art methods by a large margin. Consequently, the comprehensive ablation studies validate the individual contribution of each proposed module, while the cross-dataset evaluation confirms the strong generalization capability of our model when trained on synthetic data and directly tested on real-world scenes without fine-tuning.

Despite the promising performance, our method still exhibits certain limitations. Specifically, the computational overhead and memory consumption remain relatively high when processing light fields with ultra-high angular resolutions, which may hinder its deployment in resource-constrained, real-time applications. Additionally, the depth prediction accuracy slightly degrades in scenarios involving severe occlusions and highly reflective surfaces, where the fundamental epipolar linearity assumption is partially violated.

Finally, future research will focus on optimizing the network architecture to reduce computational complexity, potentially through lightweight feature extraction designs or adaptive angular sampling strategies. We also plan to explore the integration of self-supervised learning paradigms to alleviate the reliance on large-scale, densely annotated datasets, thereby further advancing the practical applicability of light field depth estimation in downstream tasks such as autonomous driving and robotic navigation.

Discussion

The fundamental bottleneck in light field depth estimation has long been the catastrophic "texture-copy" artifact in non-Lambertian regions. As demonstrated in our analysis, this failure is not merely a limitation of network capacity but stems from the field's entrenched reliance on the flawed view-consistency assumption \cite{wanner2014, heber2018}. By shifting the paradigm from heuristic photo-consistency to wave-vector-level angular frequency physics, our proposed physically interpretable material classification based on dual-mask and MRI-like angular frequency analysis provides a rigorous solution. Specifically, treating the $9 \times 9$ angular patch akin to an MRI k-space allows the 2D-DFT to bypass spatial heuristics and directly extract physical reflectance profiles. This mechanism effectively decouples diffuse and specular components, offering a robust material classification that traditional Epipolar Plane Image (EPI) gradient methods inherently fail to achieve \cite{li2021}.

The transition from purely data-driven feature extraction to physics-informed angular frequency analysis marks a significant milestone in 4D vision systems. While contemporary architectures excel at capturing global context, they often treat non-Lambertian reflections as noise or occlusions \cite{liu2021}. In contrast, our dual-mask mechanism explicitly quantifies pixel-level roughness and deflection vector fields, transforming anomalous reflections into informative physical priors. This physical interpretability bridges the gap between classical rendering equations and modern deep learning, enabling the system to maintain angular consistency even in the presence of complex anisotropic reflections.

Furthermore, the ternary Mask adaptive energy adjustment strategy combining multi-view geometric consistency descriptor (GCD) and SAM segmentation introduces a critical mechanism for eliminating "smooth artifacts," where high-frequency surface textures are erroneously interpreted as geometric depth structures. By generating a ternary mask—categorizing regions into rough Lambertian, smooth Lambertian, and smooth non-Lambertian—and applying adaptive energy adjustment, the network dynamically suppresses misleading texture cues while amplifying genuine structural boundaries. This strategy proves particularly vital for highly specular and transparent objects, where conventional learning-based approaches often suffer from severe depth ambiguities \cite{kirillov2023, zhang2023}. 

Ultimately, the component-aware fusion depth estimation framework based on BRDF physical parameter fitting ensures that the final depth map is not a black-box aggregation but a physically constrained synthesis \cite{blinn1977}. Unlike pure data-driven models that struggle to generalize across varying illumination conditions, our BRDF-guided fusion explicitly models light-matter interaction. This physical grounding not only enhances the fidelity of depth discontinuities but also establishes a new theoretical foundation for computational photography and virtual reality systems. Future iterations of this framework could explore integrating differentiable rendering pipelines to further close the loop between physical parameter estimation and novel view synthesis, thereby pushing the boundaries of immersive 3D reconstruction.

Limitations and Future Work

While the proposed unified dual-mask physical model significantly mitigates the texture-copying artifacts inherent in non-Lambertian regions, several limitations remain that outline promising directions for future research. 

First, the current physically interpretable material classification, which relies on wave-vector-level angular frequency analysis via 2D-DFT, primarily categorizes surfaces into diffuse, specular, and scattering profiles based on frequency energy distribution. However, this ternary simplification struggles with highly anisotropic materials or complex subsurface scattering phenomena, such as translucent skin or jade. The empirical Bidirectional Reflectance Distribution Function (BRDF) fitting \cite{blinn1977} may fail to capture the intricate multi-lobe reflectance properties of these extreme materials. Future work will explore integrating higher-order spherical harmonics or neural BRDF representations to achieve a more comprehensive physical decomposition.

Second, the adaptive energy adjustment strategy leverages the Segment Anything Model (SAM) \cite{kirillov2023} alongside the Geometric Consistency Descriptor (GCD) to generate ternary masks. While highly effective for static, well-lit scenes, SAM's zero-shot segmentation can produce boundary artifacts in the presence of severe occlusions or low-texture gradients. Furthermore, the reliance on dense $9 \times 9$ angular patches restricts the applicability of our frequency-domain analysis to high-resolution plenoptic cameras \cite{ng2005}. For sparse light field setups, angular aliasing may corrupt the high-frequency specular signatures. To address this, future iterations will investigate robust occlusion-aware masking mechanisms and super-resolution techniques in the angular domain to recover degraded frequency features.

Finally, regarding system deployment, the component-aware fusion depth estimation framework incurs substantial memory overhead due to the parallel processing of multi-scale angular attention and explicit epipolar constraints. To transition this framework from offline computational photography to real-time autonomous systems, future research will focus on architectural distillation. By incorporating efficient hierarchical vision designs \cite{liu2021} and exploring self-supervised photometric loss formulations \cite{godard2019}, we aim to reduce computational complexity without compromising physical interpretability. Additionally, embedding our wave-vector physical priors into emerging continuous neural rendering paradigms, such as Neural Radiance Fields (NeRF) \cite{ng2022}, presents an exciting avenue for jointly optimizing non-Lambertian geometry and novel view synthesis in dynamic 4D environments.

\bibliographystyle{IEEEtran}
\bibliography{references}

\end{document}
